\section{Discussion}\label{sec:discussion}

\textbf{What the NO-GO means.} Neither held-out set gives the evidence the decision rule asks
for, that a choice per document beats the best single setup.
This is not proof that a better chooser can never help: the best single setup is
selected on the same questions it is scored on, and the oracle chooses from only
half of a document's few questions.  Within the cases tested (English PDFs, two parsers, 1,250
tokens), this supports a simple default: one small-chunk setup rather than a
per-document selector.

\textbf{Fixed $k$ versus fixed budget.} RQ2 and RQ3 (the latter on tuning data)
together suggest that the best chunk size is not a property of the chunker alone.
Under a fixed number of chunks, a larger chunk shows the reader more text, so
Recall@$k$ rewards size itself; under a fixed budget of 750 or 1,250 tokens that
advantage disappears and small chunks tend to win. Comparisons at a fixed number of chunks therefore mix chunk quality with the
amount of text read; we suggest that they report the token budget, or both
measures.

\textbf{Defaults.} Tuning paid off only against LangChain's default chunk size
(post-hoc, Section~\ref{sec:rq5}). For settings like those tested, a practical check is to
compare a library's default chunk size with the token budget in use and, when a
parser without OCR is used, to detect scanned pages.
